% -----------------------------------------------------------
\section{Unfold}
% -----------------------------------------------------------
	\label{UNFOLD}
	
% % ----------------------
% \subsection{See also}
% % ----------------------

% ----------------------
\subsection{1828 American Dictionary of the English Language} % *1828
% ----------------------
	\label{UNFOLD-1828}

	\begin{compactitem}
		\item[1] To open folds; to expand; to spread out.
		\item[2] To open any thing covered or close; to lay open to view or contemplation; to disclose; to reveal; as, to \textit{unfold} one's designs; to \textit{unfold} the principles of a science.
		\item[3] To declare; to tell; to disclose.
		\item[4] To display; as, to \textit{unfold} the works of creation.
		\item[5] To release from a fold or pen; as, to \textit{unfold} sheep.
	\end{compactitem}

% % ----------------------
% \subsection{Oxford English Dictionary} % *OED
% % ----------------------
	% \label{UNFOLD-OED}

	% \begin{compactitem}
		% \item[]
	% \end{compactitem}
	
% % ----------------------
% \subsection{Restoration Lexicon} % *RL *RESTORATION LEXICON
% % ----------------------
	% \label{UNFOLD-OED}

	% \begin{compactitem}
		% \item[]
	% \end{compactitem}

% ----------------------
\subsection{Scriptural Usage} % *SCRIPTURES
% ----------------------
	\label{UNFOLD-SCRIPTURES}

	% -----
	\subsubsection{Old Testament} % *OT
	% -----
		\label{UNFOLD-OT}
	
		% --
		\paragraph{KJV}
		% --
			\label{UNFOLD-OT-KJV}
	
			\begin{tabular}{ll}
				Total occurrences in the OT & 0
			\end{tabular}
			
			% \begin{compactdesc}
				% \item[]
			% \end{compactdesc}
			
		% % --
		% \paragraph{Hebrew Bible}
		% % --
			% \label{UNFOLD-HEBREW}
		
			% %
			% \subparagraph{\texorpdfstring{\he{sample word}}{}} % paste actual hebrew text here
			% %
				% \label{PAGA} % whatever is in step bible without periods
			
				% \begin{compactdesc}
					% \item[HALOT , PDF ] \textbf{} % Use the 1994 PDF
						% \begin{compactitem}
							% \item[1]
						% \end{compactitem}
					% \item[BDB , PDF ] \textbf{}
						% \begin{compactitem}
							% \item[1]
						% \end{compactitem}
					% \item[DCH PDF ] \textbf{} % don't include the actual page number, they repeat from volume to volume
						% \begin{compactitem}
							% \item[1]
						% \end{compactitem}
				% \end{compactdesc}
				
				% \begin{compactdesc}
					% \item[Some OT uses] \textbf{}
						% \begin{compactdesc}
							% \item[]
						% \end{compactdesc}
				% \end{compactdesc}
			
		% % --
		% \paragraph{Septuagint}
		% % --
			% \label{UNFOLD-LXX}
		
			% %
			% \subparagraph{\texorpdfstring{\gr{sample word}}{}}
			% %
		
	% -----
	\subsubsection{New Testament} % *NT
	% -----
		\label{UNFOLD-NT}
	
		% --
		\paragraph{KJV}
		% --
			\label{UNFOLD-NT-KJV}		
	
			\begin{tabular}{ll}
				Total occurrences in the NT & 0
			\end{tabular}
			
			% \begin{compactdesc}
				% \item[]
			% \end{compactdesc}
			
		% % --
		% \paragraph{Greek New Testament} % *GREEK
		% % --
			% \label{UNFOLD-GREEK}
		
			% %
			% \subparagraph{\texorpdfstring{\gr{sample word}}{}}
			% %
				% \label{KATALLAGE}
			
				% \begin{compactdesc}
					% \item[BDAG , PDF ] \textbf{}
						% \begin{compactitem}
							% \item 
						% \end{compactitem}
					% \item[CGL , PDF ] \textbf{}
						% \begin{compactitem}
							% \item 
						% \end{compactitem}
					% \item[LSJ , PDF ] \textbf{}
						% \begin{compactitem}
							% \item 
						% \end{compactitem}
				% \end{compactdesc}
				
				% \begin{tabular}{ll}
					% Total occurrences in the NT & 
				% \end{tabular}
				
				% \begin{compactdesc}
					% \item[Some NT uses] \textbf{}
						% \begin{compactdesc}
							% \item[]
						% \end{compactdesc}
				% \end{compactdesc}
				
			% %
			% \subparagraph{TDNT: \texorpdfstring{\gr{sample word}}{}  (3:536--556)}
			% %
				% \label{KATALLAGE-TDNT}
		
	% -----
	\subsubsection{Book of Mormon} % *BOM
	% -----
		\label{UNFOLD-BOM}
	
		\begin{tabular}{ll}
			Total occurrences in the BOM & 10
		\end{tabular}
		
		\begin{compactdesc}
			\item[{\hyperref[1NEPHI-10-19]{1 Nephi 10:19}}]
			\item[{\hyperref[JACOB-4-18]{Jacob 4:18}}]
			\item[{\hyperref[MOSIAH-2-9]{Mosiah 2:9}}]
			\item[{\hyperref[MOSIAH-8-19]{Mosiah 8:19}}]
			\item[{\hyperref[MOSIAH-29-33]{Mosiah 29:33}}]
			\item[{\hyperref[MOSIAH-29-35]{Mosiah 29:35}}]
			\item[{\hyperref[ALMA-12-1]{Alma 12:1}}]
			\item[{\hyperref[ALMA-40-3]{Alma 40:3}}]
			\item[{\hyperref[ETHER-4-7]{Ether 4:7}}]
			\item[{\hyperref[ETHER-4-16]{Ether 4:16}}]
		\end{compactdesc}
		
	% -----
	\subsubsection{Doctrine and Covenants} % *DC
	% -----
		\label{UNFOLD-DC}
	
		\begin{tabular}{ll}
			Total occurrences in the DC & 7
		\end{tabular}
		
		Of the seven total uses of ``unfold'' in the DC, five are epistemological and four of those occur in conjunction with ``mystery''; the other two are in \hyperref[DC88-95]{DC 88:95}, which I guess you could give an epistemological interpretation of, but there are other interpretations too.
		
		For examples of ``unfolding,'' see the end of \hyperref[DC94]{DC 94}; 
		
		\begin{compactdesc}
			\item[{\hyperref[DC6-7]{DC 6:7}}] 
			\item[{\hyperref[DC10-64]{DC 10:64}}] 
			\item[{\hyperref[DC11-7]{DC 11:7}}]
			\item[{\hyperref[DC32-4]{DC 32:4}}] \phantom{.}
				\begin{compactitem}
					\item This is the one epistemological use that doesn't occur with ``mystery.''
				\end{compactitem}
			\item[{\hyperref[DC88-95]{DC 88:95}}] \textbf{(2)}
			\item[{\hyperref[DC90-14]{DC 90:14}}]
		\end{compactdesc}
		
	% -----
	\subsubsection{Pearl of Great Price} % *PGP
	% -----
		\label{UNFOLD-PGP}
	
		\begin{tabular}{ll}
			Total occurrences in the PGP & 0
		\end{tabular}
		
		% \begin{compactdesc}
			% \item[]
		% \end{compactdesc}
		
% ----------------------
\subsection{Key Phrases} % *KEYPHRASES *PHRASES
% ----------------------
	\label{UNFOLD-PHRASES}

	\begin{compactdesc}
		\item[unfold mystery] \textbf{9} | \hyperref[1NEPHI-10-19]{1 Nephi 10:19}; \hyperref[JACOB-4-18]{Jacob 4:18}; \hyperref[MOSIAH-2-9]{Mosiah 2:9}; \hyperref[MOSIAH-8-19]{Mosiah 8:19}; \hyperref[ALMA-40-3]{Alma 40:3}; \hyperref[DC6-7]{DC 6:7}; \hyperref[DC10-64]{DC 10:64}; \hyperref[DC11-7]{DC 11:7}; \hyperref[DC90-14]{DC 90:14}
			\begin{compactdesc}
				\item[Bible anchor verses] ---
					% \begin{compactdesc}
						% \item[Romans 5:5] \gr{}
					% \end{compactdesc}
				\item[Commentary] The fact that these two words occur together nine times is extremely significant, in my opinion.
			\end{compactdesc}
	\end{compactdesc}
	
% % ----------------------
% \subsection{Bible Dictionaries} % *DICTIONARIES
% % ----------------------
	% \label{UNFOLD-BD}

% ----------------------
\subsection{Restoration Usage} % *RESTORATION
% ----------------------
	\label{UNFOLD-RESTORATION}

	% % -----
	% \subsubsection{Joseph Smith} % *JS
	% % -----
		% \label{UNFOLD-JS}
	
		% \begin{compactdesc}
			% \item[{\hyperref[KEY]{Date/Source}}]
		% \end{compactdesc}
		
	% % -----
	% \subsubsection{Temple Ordinances}
	% % -----
		% \label{UNFOLD-TEMPLE}
	
		% \begin{compactdesc}
			% \item[{\hyperref[KEY]{Date/Source}}]
		% \end{compactdesc}
	
	% -----
	\subsubsection{Brigham Young} % *BY
	% -----
		\label{UNFOLD-BY}
	
		\begin{compactdesc}
			\item[{\hyperref[EMS-1865-BY-8-OCT-1865]{8 October 1865}}] I wish to say this that there is so many of the elders of Israel when they see one thing you can't get anything else in their minds you hand them a book do you see that what is on it George A Smith Brigham Young John Taylor it is a Book of Mormon they will turn [thing?] and swear there is no other side to the book there is always many sides to a thing but there is \sout{no} so many that can't see only one thing
		\end{compactdesc}
	
% % ----------------------
% \subsection{Commentary and Studies} % *COMMENTARY *STUDIES
% % ----------------------
	% \label{UNFOLD-COMMENTARY}

	% % -----
	% \subsubsection{Key Points}
	% % -----
		% \label{UNFOLD-KEYS}
	
		% \begin{compactitem}
			% \item
		% \end{compactitem}
		
	% % -----
	% \subsubsection{Sample Study}
	% % -----
		% \label{UNFOLD-study}